% =========================================================================
% Remaining sections for "Beyond Closed-Set Accuracy: Constituent-Level
% Generalization in Compound Motor Fault Diagnosis"
%
% Contents, in template order:
%   1. Abstract                  -> replace  \begin{abstract} abstract \end{abstract}
%   2. Index terms               -> replace  \begin{IEEEkeywords} keywords \end{IEEEkeywords}
%   3. Section I  Introduction   -> replace the placeholder Introduction
%   4. Section V  Conclusion     -> replace the placeholder Conclusion
%   5. Appendix A                -> \appendices after the Conclusion
%   6. Data / Code availability  -> replace the current statements
%
% Every number quoted below is taken from the current manuscript body
% (Tables 2-6 and Sections IV-A..IV-H) or, in Appendix A, from the released
% benchmark suite's multi-seed campaign CSVs. Nothing is asserted that the
% paper or the released artifacts do not contain.
% =========================================================================

% ------------------------------------------------------------------ 1. ABSTRACT
\begin{abstract}
Reported accuracies for data-driven motor fault diagnosis are frequently near
ceiling, yet they rarely measure the ability that industrial deployment
actually requires: recognizing the elementary fault mechanisms of a compound
fault whose exact combination was never observed during model development.
This paper evaluates that ability on a public multimodal dataset of a 2.2-kW
three-phase induction motor containing 24 diagnostic states, including nine
compound faults that pair a mechanical bearing defect with a second
mechanical, electrical, or electromagnetic mechanism, under twelve speed/load
conditions. The problem is formulated at the constituent level: nine
elementary mechanisms are detected independently from physics-guided,
mechanism-specific vibration- and current-domain features, so unseen
combinations remain expressible. A crossed protocol separates three levels of
compound-context exposure (recombination, one-sided, and two-sided
context zero-shot) from operating-condition novelty, with run-level,
severity-aware splits and all calibration confined to training data. Although
closed-set 24-class accuracy reaches 0.913, exact recovery of unseen
constituent sets remains between 0.083 and 0.157 for the locked linear probe,
and every metric is reported against input-independent prevalence references
that expose partial-recovery scores as otherwise misleading. Compound-context
deprivation is measurably harmful chiefly when the operating condition is
also unseen (micro-$F_1$ interaction $+0.040$, 95\% CI $[0.009, 0.073]$), and
cross-partner transfer is strongly directional, including one statistically
supported score-orientation reversal whose physical attribution is
deliberately left open given demonstrated recording-session separability.
Threshold-objective and classifier-family analyses show that exact
decomposition, not constituent recall, is the binding constraint.
\end{abstract}

% ---------------------------------------------------------------- 2. KEYWORDS
\begin{IEEEkeywords}
compound faults, condition monitoring, constituent-level diagnosis, data
leakage, electrical machines, evaluation protocols, fault diagnosis,
induction motors, multi-label classification, zero-shot learning
\end{IEEEkeywords}

% ------------------------------------------------------------ 3. INTRODUCTION
\section{Introduction}
\label{sec:intro}

\PARstart{I}{nduction} motors underpin industrial drivetrains, and their
unplanned failures propagate directly into downtime and safety risk.
Data-driven condition monitoring has consequently produced a large literature
whose reported accuracies now cluster near ceiling on standard laboratory
benchmarks. Deployment experience is harder to reconcile with those numbers:
a model commissioned at one operating point degrades at another, and faults
in service rarely arrive one at a time. A bearing defect that develops
alongside a winding short circuit, broken rotor bars, or eccentricity
confronts the monitoring system with a \emph{compound} state whose exact
combination was, in general, never present in the training history.

Conventional closed-set formulations cannot even express this situation: a
previously unseen combination $A{+}B$ constitutes an unseen class. The
compound-fault literature therefore moved toward decoupling formulations that
predict elementary mechanisms rather than combinations, and, more recently,
toward zero-shot settings in which compound states must be recognized from
single-fault training data alone. As Section~\ref{sec:related} reviews,
however, that literature has evaluated almost exclusively \emph{homogeneous}
compounds---bearing--bearing or bearing--gear pairings whose constituents are
all observable in the same vibration channel---and its reported accuracies of
roughly 75--87\% have been obtained under evaluation practices that a
parallel line of work on data leakage has shown to be fragile: sample-level
splits, steady operating conditions, and metrics reported without reference
to what label prevalence alone would achieve.

This paper asks a deliberately harder question on a public dataset that makes
it answerable. The MCC5-THU motor dataset \cite{chen2026} provides 288
recordings of a 2.2-kW three-phase induction motor across 24 diagnostic
states under twelve speed/load conditions, with synchronized triaxial
vibration, three-phase currents, torque, and a key-phase tachometer. Its nine
compound faults pair a mechanical bearing defect with a second
mechanism---inner or outer race, broken rotor bars, static or dynamic
eccentricity, or a stator winding short circuit---so that, in eight of nine
combinations, the two constituents deposit their primary evidence in
\emph{different} sensing modalities. To our knowledge, neither this
cross-modal compound regime nor this dataset has previously received a
constituent-level evaluation.

The study is designed so that every claim is bounded by an explicit control.
The formulation is constituent-level: nine elementary mechanisms are detected
independently by regularized linear probes operating on curated,
mechanism-specific physical-signature families derived from angle-resampled
order spectra, envelope-order descriptors, and phase-current sequence
measures. A crossed protocol varies compound-context exposure over three
regimes---recombination, one-sided, and two-sided compound-context
zero-shot---and crosses each with seen and unseen operating conditions, under
run-level, severity-aware splits with all preprocessing, calibration, and
threshold selection confined to training data. Every reported metric is
accompanied by input-independent constant-predictor references, so that
prevalence-driven scores cannot masquerade as diagnostic skill; oracle
detectability analyses are paired with recording-sensitivity and same-date
controls, so that separability is not silently attributed to fault
mechanisms; and paired run-clustered bootstrap inference with multiplicity
correction distinguishes supported effects from noise.

The contributions are as follows.
\begin{enumerate}
  \item \textbf{A constituent-level formulation and crossed
  composition--context--condition protocol} for compound motor faults whose
  constituents span mechanical and electrical modalities, separating novel
  composition, compound-context deprivation, and operating-condition novelty
  in one design (Sections~\ref{sec:methods}--\ref{sec:results}).
  \item \textbf{Prior-aware evaluation}: all constituent metrics are reported
  against input-independent prevalence references, which expose several
  apparently favorable partial-recovery metrics as at or below what a
  constant predictor achieves, while establishing that constituent recall and
  all-constituent recovery genuinely exceed those floors.
  \item \textbf{A controlled analysis of constituent transfer}: within-pair
  oracle detectability, cross-partner transfer, physical-signature
  preservation, and a statistically supported score-orientation reversal,
  each bounded by recording-sensitivity and same-date controls that keep
  session effects from being read as physics.
  \item \textbf{Sensitivity and robustness analyses} over the decision-threshold
  objective, single-fault severity merging, and four classifier families,
  showing that the qualitative generalization pattern is stable while
  absolute exact-match decomposition is an operating-point and model
  property, not an intrinsic bound of the protocol.
  \item \textbf{Documentation of dataset-release properties that materially
  affect protocol design}---recording-session separability, severity-variant
  leakage paths, and label abbreviations---together with reference
  operating-condition experiments (Appendix~\ref{app:condition}) that
  motivate the crossed design.
\end{enumerate}

The principal findings are that closed-set discriminability (accuracy 0.913)
coexists with exact constituent-set recovery of 0.083--0.157 on unseen
compounds for the locked probe; that recovery of at least one constituent is
nearly universal while exact decomposition fails, so the practical failure
mode is incomplete decomposition with over-diagnosis rather than silence;
that compound-context deprivation becomes measurably harmful specifically
under simultaneous operating-condition novelty; and that constituent evidence
is often, but directionally and heterogeneously, transferable across
compound partners.

The remainder of the paper is organized as follows.
Section~\ref{sec:related} reviews related work and positions the study.
Section~\ref{sec:methods} details the constituent representation, feature
construction, detection, protocol, metrics, and statistical analysis.
Section~\ref{sec:results} reports results. Section~\ref{sec:conclusion}
concludes. Appendix~\ref{app:condition} reports the operating-condition
reference experiments.

% ------------------------------------------------------------- 4. CONCLUSION
\section{Conclusion}
\label{sec:conclusion}

This paper evaluated whether the elementary mechanisms of a compound motor
fault can be recovered when their exact combination, their compound context,
and the operating condition are varied independently of the training history.
On a public multimodal induction-motor dataset whose compound faults span
mechanical and electrical modalities, a locked linear constituent probe over
mechanism-specific physical features showed that strong closed-set
discriminability (24-class accuracy 0.913, 95\% CI $[0.885, 0.938]$) coexists
with exact constituent-set recovery of only 0.083--0.157 on unseen compounds.
The two results are not in tension; they measure different abilities, and the
gap between them is the paper's subject.

Four conclusions follow from the controlled analyses. First, the failure
mode is incomplete decomposition rather than loss of information: at least
one true constituent was recovered in more than 92\% of compound runs, and
constituent recall (0.630--0.694) and all-constituent recovery under
recombination (up to 0.444) exceed input-independent prevalence references
(0.556 and 0.111, respectively), while several other partial-recovery
metrics do not---a distinction visible only because every metric was
reported against its constant-predictor floor. Second, compound-context
deprivation is measurably harmful chiefly when the operating condition is
simultaneously unseen: the recombination-versus-two-sided penalty under
condition shift was $+0.040$ micro-$F_1$ (95\% CI $[0.009, 0.073]$),
$+0.051$ constituent recall, and $+0.083$ all-constituent recovery, whereas
the corresponding seen-condition contrasts were not resolved. Third,
constituent evidence transfers across compound partners often but
directionally: within-pair detectability was near-perfect throughout, yet
cross-partner transfer ranged from perfect to a statistically supported
score-orientation reversal (Holm-adjusted $p \le 0.014$)---and the
demonstrated separability of distinct recordings of the same fault state
requires that reversal to be interpreted as context-dependent score
transfer rather than as proven physical signature inversion. Fourth,
absolute decomposition quality is an operating-point and model property:
precision-weighted thresholding raised exact-match recovery to 0.157--0.287
at reduced recall, and nonlinear classifier families reduced false additions
from ${\sim}0.68$ to ${\sim}0.26$ per run while the qualitative
generalization pattern persisted across all four families.

For practice, these results caution against reading closed-set accuracy as
deployment readiness, argue for calibrating constituent detectors with
explicit precision--recall cost structure, and indicate that commissioning
data should sample compound contexts and operating conditions jointly rather
than relying on either alone. The reference experiments of
Appendix~\ref{app:condition} reinforce the latter point: on this dataset the
standard leave-one-condition-out test is nearly saturated, while
single-condition training collapses accuracy, so robustness claims should be
evaluated in the extrapolating direction.

Several directions remain open. The dataset provides no replicated
same-state recordings that would separate fault, assembly, sensor-mounting,
and session effects, so the transfer analyses stop at carefully bounded
attribution; replicated acquisitions would convert those bounds into causal
statements. Representation or normalization strategies that promote
constituent-score consistency across compound partners are identified as
hypotheses, not established remedies. Finally, the constituent-level
protocol, prevalence references, and controls used here are directly
portable to the companion gearbox release and to other multimodal datasets,
where cross-dataset replication of the context-by-condition interaction
would test its generality.

% ------------------------------------------------------------- 5. APPENDIX A
\appendices
\section{Operating-Condition Reference Experiments}
\label{app:condition}

Section~\ref{sec:related} states that the leave-one-condition-out protocol,
the standard robustness test in the multi-condition literature, is nearly
saturated on this dataset, whereas the reverse direction is not. This
appendix reports the experiments behind that statement. They use a
deliberately generic pipeline---distinct from, and simpler than, the
constituent-level methodology of Section~\ref{sec:methods}---so that the
measured gaps reflect protocol difficulty rather than modeling choices, and
they are included with the released code.

Windows of 0.64\,s (8192 samples, non-overlapping) were described by 108
time- and frequency-domain features (eight temporal statistics and ten
spectral descriptors per channel over the three vibration and three current
channels), classified into the original 24 diagnostic states by a
300-tree random forest, with three seeds per configuration. Five split
protocols were compared. \emph{Random windows} splits windows irrespective
of recording (a deliberately leakage-prone reference); \emph{in-condition}
trains on the first 60\% of every recording and tests on the final 30\% with
a 10\% guard gap; \emph{unknown condition} holds out each of the twelve
operating conditions in turn; \emph{cross-profile} trains on one excitation
profile (constant-torque variable-speed or constant-speed variable-torque)
and tests on the other; and \emph{single source} inverts the unknown-condition
protocol, training on one condition and testing on the remaining eleven. A
sixth split trains on quasi-stationary segments and tests on speed/load
ramps, with stationarity labeled from the key-phase-derived speed and the
torque channel; the torque channel is used only for this protocol
construction and never as a classifier input.

\begin{table}[!t]
\caption{Reference protocols, 24-class accuracy (mean $\pm$ std over folds
and three seeds). The same features and classifier are used throughout; only
the split changes.}
\label{tab:condref}
\centering
\footnotesize
\setlength{\tabcolsep}{5pt}
\begin{tabular}{l c c}
\hline
Protocol & Folds & Accuracy \\
\hline
Random windows (leakage-prone reference) & 1 & $0.975 \pm 0.000$ \\
In-condition (guard-gap temporal split)  & 1 & $0.965 \pm 0.000$ \\
Unknown condition (leave-one-out)        & 12 & $0.939 \pm 0.024$ \\
Cross-profile                            & 2 & $0.839 \pm 0.115$ \\
Steady $\rightarrow$ transitional segments & 1 & $0.766 \pm 0.001$ \\
Single source (train on one condition)   & 12 & $0.333 \pm 0.104$ \\
\hline
\end{tabular}
\end{table}

Table~\ref{tab:condref} summarizes the outcome. Holding one condition out of
twelve costs only ${\sim}2.6$ accuracy points relative to the in-condition
reference, because the eleven retained conditions bracket the held-out one
in speed and load and the model interpolates. Training on a single condition
and extrapolating to the other eleven costs over 60 points; paired over the
twelve condition folds, the difference between the two directions is 0.606
(paired $t$-test $p = 7.7\times10^{-10}$; Wilcoxon $p = 4.9\times10^{-4}$).
Transfer between excitation profiles is likewise asymmetric at equal
training size: training under variable speed generalizes to variable load
(accuracy 0.944), whereas the reverse loses ${\sim}21$ points (0.734; paired
over seeds, $p = 3.3\times10^{-5}$). Finally, appending 36 speed-invariant
envelope-order descriptors to the feature set leaves the interpolating
protocols unchanged ($+0.001$ on unknown-condition, paired $p = 0.84$) but
significantly improves the extrapolating one ($+0.047$ on single source,
paired $p = 8.4\times10^{-4}$), consistent with their construction.

Two implications carry into the main study. Methodologically, robustness to
operating conditions should be claimed in the extrapolating direction, which
motivates crossing the composition and context regimes with unseen-condition
evaluation rather than relying on leave-one-condition-out alone.
Practically, excitation diversity during commissioning matters more than
data volume: a training condition in which speed varies transfers to load
variation, but not conversely.

% ------------------------------------------- 6. DATA AND CODE AVAILABILITY
% Replace the current DATA AVAILABILITY / CODE AVAILABILITY statements with:

\section*{Data Availability}
The dataset analyzed in this study is publicly available from Mendeley Data
at \mbox{https://doi.org/10.17632/6s3dggj9mw.1} (CC~BY~4.0)
\cite{mcc5motor_data}, with its accompanying descriptor article
\cite{chen2026}. Acquisition-timestamp metadata used for the
recording-sensitivity controls were recovered from the released file names
and are included with the analysis code.

\section*{Code Availability}
The complete analysis code---including the constituent-level pipeline, the
crossed evaluation protocol, the constant-predictor references, the
recording-sensitivity controls, and the reference experiments of
Appendix~\ref{app:condition}---is available from the corresponding author
upon reasonable request.
% Recommended alternative (stronger for a paper whose subject is evaluation
% reliability): deposit the repository at Zenodo and cite its DOI here, e.g.
% "is openly available at https://doi.org/10.5281/zenodo.XXXXXXX".
